\chapter{Fusion lattices and their higher centre}
\label{chap:fc-fusion}

A rational exchange operator compares two tensor representations wherever
its denominator is nonzero. At a pole, that comparison can still determine
an algebra, but only if we retain the modules on both sides of the exchange.
The smallest example already distinguishes three objects: the algebra at
the pole, its action on the first limiting module, and its action on the
second. Each action loses an arrow that the algebra retains.

We shall compute this algebra directly from the tensor operators. Its
relative centre will have one class in every positive even degree supported
at fusion. More decisively, its multiplication can be replaced by the
cohomology multiplication, whereas its insertion bracket cannot be replaced
by the cohomology bracket over the same coefficient ring. The equation
\[
 d\theta=fu,\qquad [\theta,u]=-u
\]
will express both facts. Here \(f\) is the equation of the fusion divisor;
the degrees of \(\theta\) and \(u\) will be cochain degrees, not spatial
dimensions or Fourier indices.

\section{Two modules at a spectral resonance}

Fix an integer \(d\geq2\), let \(V=\mathbb C^d\), and write \(e_i\),
\(0\leq i<d\), for its basis. The matrix unit \(E_{ij}\) sends \(e_j\)
to \(e_i\) and kills the other basis vectors. On \(V\otimes V\), let
\(P(v\otimes w)=w\otimes v\). The idempotents
\[
 P_+=(1+P)/2,\qquad P_-=(1-P)/2
\]
project onto the symmetric and alternating tensors. Their dimensions are
\(p=d(d+1)/2\) and \(q=d(d-1)/2\).

Fix \(\hbar\in\mathbb C^\times\). Put \(R_0=\mathbb C[s]\) and
\(K=\mathbb C(s)\). A lattice in a finite-dimensional \(K\)-vector
space means a free \(R_0\)-submodule that spans it over \(K\).
Consider the exchange operator and the two lattices
\begin{equation}\label{eq:fc-lattices}
 \mathsf R(s)=\frac{sI-\hbar P}{s-\hbar}
   =P_++\frac{s+\hbar}{s-\hbar}P_-,\qquad
 M=R_0\otimes(V\otimes V),\qquad L=\mathsf R(s)^{-1}M.
\end{equation}
Thus \(L=M_+\oplus\frac{s-\hbar}{s+\hbar}M_-\), where
\(M_\pm=P_\pm M\). Both lattices are considered in the same rational
vector space. The simultaneous endomorphism algebra consists of the
\(K\)-linear maps preserving each lattice.

We fix the representation conventions before calculating this algebra.
Let \(\mathsf Y_\hbar(d)\) be generated by the coefficients of series
\(T_{ij}(w)=\delta_{ij}+\sum_{n\geq1}t_{ij}^{(n)}w^{-n}\), subject to
\begin{equation}\label{eq:fc-rtt}
 (w-v)[T_{ij}(w),T_{kl}(v)]
 =\hbar\bigl(T_{kj}(w)T_{il}(v)-T_{kj}(v)T_{il}(w)\bigr).
\end{equation}
This is an algebraic relation on every Laurent coefficient. There is no
analytic convergence condition on the series. The evaluation action at a
scalar \(z\) sends \(T_{ij}(w)\) to
\(\delta_{ij}+\hbar E_{ij}/(w-z)\). The coproduct is matrix
multiplication, \(\Delta T_{ij}=\sum_aT_{ia}\otimes T_{aj}\).

For completeness, the evaluation relation follows from
\([E_{ij},E_{kl}]=\delta_{jk}E_{il}-\delta_{il}E_{kj}\): substituting
the evaluation formula in \eqref{eq:fc-rtt} cancels its quadratic terms
and leaves that matrix identity. For the coproduct, write the relation as
\((1-\hbar P/(w-v))T_1(w)T_2(v)
=T_2(v)T_1(w)(1-\hbar P/(w-v))\), where the subscripts refer to two
auxiliary matrix spaces. In a product of two evaluation matrices, the
entries acting on different quantum factors commute. Moving those entries
past each other and using this equation once in each factor proves the
same equation for the product. Associativity of matrix multiplication
proves coassociativity.

Use evaluation parameters \(0,s\), in that order. Write
\(A_{ij}=E_{ij}\otimes1\), \(B_{ij}=1\otimes E_{ij}\), and
\(J_{ij}=A_{ij}+B_{ij}\). The first two coefficients are
\begin{equation}\label{eq:fc-coefficients}
 T^{(1)}_{ij}=\hbar J_{ij},\qquad
 T^{(2)}_{ij}=\hbar sB_{ij}+\hbar^2 A_{ij}P.
\end{equation}
Indeed \(\sum_aE_{ia}\otimes E_{aj}=A_{ij}P\), as is seen by applying
both sides to \(e_k\otimes e_l\). In particular
\(K_{ij}:=T^{(2)}_{ij}-sT^{(1)}_{ij}
=-\hbar sA_{ij}+\hbar^2A_{ij}P\).

\begin{theorem}\label{thm:fc-image}
The \(R_0\)-algebra generated by the tensor evaluation representation is
exactly
\begin{equation}\label{eq:fc-global-order}
 \Lambda_d=
 \begin{pmatrix}
  \operatorname{Mat}_p(R_0)&(s+\hbar)\operatorname{Mat}_{p\times q}(R_0)\\
  (s-\hbar)\operatorname{Mat}_{q\times p}(R_0)&\operatorname{Mat}_q(R_0)
 \end{pmatrix}.
\end{equation}
It is the simultaneous endomorphism algebra of \(M,L\).
The first two coefficients already generate it.
\end{theorem}
\begin{proof}
An integral block matrix preserves \(L\) precisely when its upper-right
block is divisible by \(s+\hbar\) and its lower-left block by
\(s-\hbar\). To see this, conjugate by
\(\operatorname{diag}(1,(s-\hbar)/(s+\hbar))\) and require every entry
to be polynomial. The two linear polynomials are coprime because
\(2\hbar\ne0\).

The operator \(\mathsf R(s)\) intertwines the coproduct and opposite
coproduct. This can also be verified directly here. Clear the evaluation
denominators and replace \(A_{ij}P\) by \(B_{ij}P\) for the opposite
coproduct. Use \(PA_{ij}=B_{ij}P\) and \(P^2=1\). For \(I+cP\)
the difference of the two sides is
\((\hbar^2+\hbar cs)(A_{ij}-B_{ij})P\); it vanishes at
\(c=-\hbar/s\). Multiplying by \(s/(s-\hbar)\) gives
\eqref{eq:fc-lattices}, including its value by rational continuation
where this intermediate expression used \(s^{-1}\). Both tensor
representations have polynomial coefficients in \(s\). If \(v\in L\),
then \(\mathsf R(s)v\in M\), so the intertwining identity implies
\(\mathsf R(s)T_{ij}^{(n)}v\in M\). Thus every coefficient preserves
both lattices. This proves one inclusion.

For the reverse inclusion, the first coefficients give every \(J_{ij}\).
They give the operator
\[
 \sum_{i,j}J_{ij}J_{ji}=2dI+2P.
\]
The two single-factor sums are \(dI\) each, and the two cross sums are
\(P\) each. Hence \(P_+\) and \(P_-\) belong to the image.
Choose the diagonal one-factor operator with entries \(3^i\).
Its diagonal tensor action has eigenvalues \(3^i+3^j\).
Their ternary expansions distinguish all unordered pairs \(i\leq j\).
Polynomial interpolation in this operator, followed by \(P_+\) or
\(P_-\), therefore projects onto each symmetric or alternating weight
line. All denominators used in interpolation are nonzero complex numbers.

The operators \(J_{ij}\) connect every pair of weight lines in either
channel by a sequence of nonzero maps. In the symmetric channel, replace
one tensor index at a time; this connects every unordered pair to a
diagonal pair. In the alternating channel, replace one index by an index
not already present. The graph of two-element subsets is connected when
\(d>2\), and has a single vertex when \(d=2\).
Sandwich a nonzero connecting map between its two weight projections,
divide by its nonzero scalar, and compose along a path. This constructs
every matrix unit on each channel independently.

It remains to connect the channels. Let
\(v=e_0\otimes e_1+e_1\otimes e_0\) and
\(w=e_0\otimes e_1-e_1\otimes e_0\). In the ordered basis \((v,w)\),
the projected operator \(K_{00}\) is
\[
 -\frac{\hbar}{2}
 \begin{pmatrix}s-\hbar&s+\hbar\\s-\hbar&s+\hbar\end{pmatrix}.
\]
The term \(A_{00}\) selects \(e_0\otimes e_1\), and
\(A_{00}P\) sends \(e_1\otimes e_0\) to that vector, which proves the
matrix. Projecting its off-diagonal entries and dividing by
\(-\hbar/2\) gives \((s+\hbar)E_{vw}\) and
\((s-\hbar)E_{wv}\). The channel matrix units now give every entry of
\eqref{eq:fc-global-order}. These entries form an \(R_0\)-basis, proving
the other inclusion without inverting either resonance.
\end{proof}

The fusion divisor is \(s=\hbar\) or \(s=-\hbar\). It is distinct
from the coincident spectral values \(s=0\) at the fixed nonzero
\(\hbar\) used here. Taking \(\hbar=0\) would make the evaluation
coefficients vanish and invalidate the generation argument.

\section{The algebra and its two limiting representations}

Set \(t=s-\hbar\) and localize to
\(R=\mathbb C[t,(t+2\hbar)^{-1}]\). The upper-right ideal in
\eqref{eq:fc-global-order} is now the whole ring. Choose one basis line
in each channel, and let \(e\) project onto their sum. The corner
\(e\Lambda_de\) has the four-element basis
\[
 e_0,e_1,a,b,\qquad a=E_{01},\quad b=tE_{10},\quad
 ab=te_0,\quad ba=te_1.
\]
The idempotent \(e\) is full: this means \(\Lambda_de\Lambda_d=\Lambda_d\).
Indeed every diagonal matrix unit can be written as a product of two
channel matrix units through its selected line, and their sum is \(1\).

Here is the module equivalence implied by this identity. For a ring
\(A\), an idempotent \(e\), and an expression
\(1=\sum_i x_i e y_i\), multiplication induces isomorphisms
\[
 Ae\otimes_{eAe}eA\longrightarrow A,\qquad
 eA\otimes_A Ae\longrightarrow eAe.
\]
The first inverse sends \(a\) to \(\sum_i ax_ie\otimes ey_i\).
For \(ae\otimes eb\), balancing gives
\(\sum_i aebx_ie\otimes ey_i=ae\otimes eb\sum_i x_iey_i\),
which proves the other composite. The second inverse sends \(c\) to
\(e\otimes c\); balancing proves its two composites. Consequently
\(N\mapsto eN\) and \(U\mapsto Ae\otimes_{eAe}U\) are inverse
functors on left modules. They are exact, since inverse additive
equivalences preserve kernels and cokernels, and apply term by term to
complexes. This is the Morita equivalence used below.
It also identifies \(A\) with \(\operatorname{End}_{eAe}(Ae)\).
For an endomorphism \(F\) of the right \(eAe\)-module \(Ae\), set
\(a_F=\sum_i F(x_ie)ey_i\). Then
\(a_Fbe=\sum_iF(x_ie)(ey_ibe)=F(be)\).
Conversely an element killing \(Ae\) kills
\(\sum_i x_iey_i=1\), so is zero. This proves the endomorphism
identification by explicit inverse maps.

We shall compute the corner over a general commutative ring \(R\).
For \(f\in R\), let \(A_f\) be free on \(e_0,e_1,a,b\), with
\begin{gather}\label{eq:fc-rankfour}
 e_i e_j=\delta_{ij}e_i,\quad 1=e_0+e_1,\quad
 a=e_0ae_1,\quad b=e_1be_0,\\
 ab=fe_0,\qquad ba=fe_1.\notag
\end{gather}
All products with incompatible endpoints are zero. The only three-arrow
associativity checks are \((ab)a=a(ba)=fa\) and its counterpart for
\(b\). Products containing a vertex follow from the endpoint rules.
Thus these formulas define an associative algebra on a free rank-four
module, even when \(f\) is a zero divisor.

For the local fusion parameter \(f=t\), the selected parts of \(M,L\)
give representations
\begin{equation}\label{eq:fc-probes}
 M:\quad a=\begin{pmatrix}0&1\\0&0\end{pmatrix},\quad
 b=\begin{pmatrix}0&0\\t&0\end{pmatrix};
 \qquad
 L:\quad a=\begin{pmatrix}0&t\\0&0\end{pmatrix},\quad
 b=\begin{pmatrix}0&0\\1&0\end{pmatrix}.
\end{equation}
The second uses the lattice basis \((v,tw)\). At \(t=0\), the first
representation kills \(b\), and the second kills \(a\). In the algebra
\(A_t/tA_t\), both are nonzero basis vectors and \(ab=ba=0\).
Specializing an embedding into matrices can therefore destroy its
injectivity. This example is the reason to retain the algebra itself.

Over the unlocalized ring \(R_0\), the selected corner has arrows
\((s+\hbar)E_{01}\), \((s-\hbar)E_{10}\), and is
\(A_{s^2-\hbar^2}\). The same calculations therefore describe both
resonances with their actual equation.

\section{Relations and the complete sequence of their primitives}

Let \(E=Re_0\oplus Re_1\). The projection killing \(a,b\) is not
an algebra augmentation \(A_f\to E\) unless \(f=0\): it sends \(ab\)
to zero while fixing \(fe_0\). To resolve the relations without
inventing that augmentation, retain a generator for every alternating
word in \(a,b\).

Write \(q_w\) for the generator associated to a nonempty alternating
word \(w\), with the same left and right endpoints and degree
\(1-|w|\). Let \(Q_f\) be the free graded path algebra on these
generators over \(E\). Thus its basis consists of finite composable
products of generators; noncomposable products are zero. Define a
degree-one derivation by
\begin{equation}\label{eq:fc-cobar}
 dq_w=\sum_{j=1}^{|w|-1}(-1)^{j+1}
             q_{w_{\leq j}}q_{w_{>j}}
       -\mathbf1_{|w|=2}fe_{\mathrm{left}(w)}.
\end{equation}
The Leibniz rule is \(d(xy)=dx\,y+(-1)^{|x|}x\,dy\).
In particular \(dq_{ab}=q_aq_b-fe_0\) and
\(dq_{aba}=q_aq_{ba}-q_{ab}q_a\). The first is the relation;
the second retains the compatibility between its two products with
an arrow.

\begin{proposition}\label{prop:fc-cobar}
The differential in \eqref{eq:fc-cobar} squares to zero. The map
\(Q_f\to A_f\) sending \(q_a,q_b\) to \(a,b\) and all longer
generators to zero is a quasi-isomorphism for every \(R,f\).
\end{proposition}
\begin{proof}
In \(d^2q_w\), each division into three consecutive nonempty blocks
appears in two orders, with opposite signs from the displayed formula
and the Leibniz rule. The remaining terms remove a two-letter end block
by its scalar curvature. Removing the first or the last two letters
leaves the same shorter alternating word. The two coefficients are
opposite multiples of \(f\). This also covers length three; length two
has zero square immediately. Hence \(d^2=0\).

First set the scalar term to zero. At a fixed underlying alternating
word of length \(n\), a basis element of \(Q_0\) specifies which of
the \(n-1\) gaps divide two generators. Its differential changes an
undivided gap to a divided gap with coefficient \(\pm1\).
The complex is a tensor product of \(n-1\) two-term identity complexes,
with the tensor differential signs. For \(n>1\), contraction at the
first gap contracts it. For \(n=1\), just the arrow survives. The
length-zero part is \(E\). This proves the assertion at \(f=0\).

For general \(f\), filter by total underlying arrow length, giving
\(A_f\) length zero on vertices and length one on arrows. Products
of arrows in \(A_f\) can lower length. The associated graded target
is \(A_0\), and the associated graded source is \(Q_0\).
The curvature term in \eqref{eq:fc-cobar} lowers length by two.
In the cone of the proposed map, every cycle is a finite sum and has
a greatest length. The calculation at \(f=0\) supplies a primitive for
its highest-length part. Subtract its differential and repeat at a
strictly smaller length. This terminates and contracts every cone
cycle. It proves the quasi-isomorphism without a convergence assumption.
\end{proof}

One may record the same curvature before taking cobar. The graded path
algebra on two degree-one arrows \(\alpha,\beta\) has differential
zero and curvature \(f(\alpha\beta+\beta\alpha)\).
The curvature is central and closed. Although its commutator is zero,
it must be retained: its two scalar components are the last terms in
\eqref{eq:fc-cobar}. All sums here are algebraic finite sums. Allowing
infinite path series would define a different completion.

\section{The diagonal and all insertion operations}

The relative centre over \(R\) is the complex of derived bimodule
endomorphisms of \(A_f\). We now compute it through a projective
resolution, keeping its actual composition operations. Put
\(J=A_f/E=Ra\oplus Rb\), an \(E\)-bimodule. The normalized bar
resolution has terms
\[
 P_n=A_f\otimes_E J^{\otimes_E n}\otimes_E A_f,\qquad n\geq0.
\]
Its differential multiplies adjacent factors with alternating signs;
an interior product is projected to \(J\). The augmentation multiplies
the two factors of \(P_0\). Exactness can be checked before applying
bimodule Hom: prepend \(1\) and the \(J\)-part of the first factor
to define an extra degeneracy. Its anticommutator with the differential
is the identity on the augmented complex. The omitted \(E\)-part
moves across the tensor sign by balancing, so this calculation does
not require an algebra augmentation \(A_f\to E\).

There are two alternating words in each positive length. Each term
is a sum of modules \(A_fe_i\otimes_R e_jA_f\), projective over
\(A_f\otimes_R A_f^{\mathrm{op}}\). Thus the cochains
\(C^n=\operatorname{Hom}_{E-E}(J^{\otimes_E n},A_f)\) compute the
relative centre. In degree zero this means the centralizer of \(E\).
The separability element \(e_0\otimes e_0+e_1\otimes e_1\) permits
the same normalization of the ordinary relative bar. Equivalently,
both are projective resolutions of the same bimodule. Extending an
\(E\)-balanced normalized cochain to all inputs by projection to
\(J\) gives its usual Hochschild representative; balancing makes its
differential, cup product and insertions the same operations.

In arity \(2m\), write a cochain as \((x,y)\): its values on the
words starting in \(a,b\) are \(xe_0,ye_1\). In arity \(2m+1\),
write \((r,s)\) for values \(ra,sb\). For an \(n\)-input cochain,
the Hochschild differential has endpoint terms
\(x_1F(x_2,\ldots)+(-1)^{n+1}F(\ldots,x_n)x_{n+1}\)
and interior terms
\(\sum_{i=1}^n(-1)^iF(\ldots,x_ix_{i+1},\ldots)\).
The interior terms vanish on arrow words because their products lie
in \(E\). Consequently
\begin{equation}\label{eq:fc-periodic}
 d_{2m}(x,y)=(y-x,x-y),\qquad
 d_{2m+1}(r,s)=f(r+s,r+s).
\end{equation}

Define \(u^m=(1,1)\) in degree \(2m\) and
\(\theta u^m=(0,1)\) in degree \(2m+1\). Their span is the
differential graded cup algebra
\begin{equation}\label{eq:fc-small}
 \mathcal C_f=R[u]\otimes\Lambda(\theta),\qquad
 |u|=2,\quad |\theta|=1,\quad d\theta=fu.
\end{equation}
Here \(\Lambda(\theta)\) means the algebra with \(\theta^2=0\).
For instance, the cup square vanishes because the only odd word where
\(\theta\) is nonzero starts with \(b\); the following odd block
then starts with \(a\). The inclusion \(i:\mathcal C_f\to C\)
has the explicit chain projection and homotopy
\begin{gather}\label{eq:fc-contraction}
 p_{2m}(x,y)=xu^m,\qquad p_{2m+1}(r,s)=(r+s)\theta u^m,\\
 h_{2m+1}(r,s)=(0,r),\qquad h_{2m}=0.\notag
\end{gather}
Substitution in \eqref{eq:fc-periodic} gives
\(pi=1\), \(dh+hd=1-ip\), and \(hi=ph=h^2=0\).

A brace \(F\{G_1,\ldots,G_l\}\) inserts the inner cochains into
distinct ordered slots of \(F\), summing all choices. Inserting an
\(n\)-input cochain after \(j\) preceding inputs contributes sign
\((-1)^{(n-1)j}\); tensor signs supply the corresponding convention
when more insertions are made. The empty brace is \(F\) itself.
These definitions include cochains with zero inputs.

\begin{theorem}\label{thm:fc-braces}
The inclusion of \(\mathcal C_f\) preserves the differential, cup
product, and every brace, and is a quasi-isomorphism. For \(l\geq1\),
its nonzero braces are
\begin{align}
 u^m\{\theta u^{n_1},\ldots,\theta u^{n_l}\}
 &=\binom ml u^{m+\sum_i n_i},\label{eq:fc-braces}\\
 \theta u^m\{\theta u^{n_1},\ldots,\theta u^{n_l}\}
 &=\binom{m+1}{l}\theta u^{m+\sum_i n_i}.\notag
\end{align}
A brace containing an even inner cochain is zero, including an
insertion of the degree-zero unit.
\end{theorem}
\begin{proof}
An even inner cochain returns a vertex scalar. A normalized outer
cochain vanishes on such an input. An odd inner cochain returns \(b\)
and is zero on a segment starting with \(a\). It therefore occupies
a \(b\)-slot of the collapsed outer word. There are \(m\) such
slots in an even word of length \(2m\), and \(m+1\) in the nonzero
odd word of length \(2m+1\). Choose \(l\) of these slots in order.
Replacing the chosen slot by a segment of length \(2n_i+1\) adds
\(2n_i\) inputs without changing any subsequent parity. All insertion
signs are positive. This gives exactly \eqref{eq:fc-braces}.
The operations are restrictions of actual cochain operations, and
\eqref{eq:fc-contraction} proves the quasi-isomorphism.
\end{proof}

\subsection{Retaining the matrix multiplicities}

The computation also gives the relative centre of the actual Yangian
image, with all its channel multiplicities. In \(\Lambda_d\), put
\(H=\operatorname{Mat}_p(R)\oplus\operatorname{Mat}_q(R)\).
Normalize the diagonal bar over \(H\), rather than over its individual
diagonal entries. Each matrix block is separable over \(R\):
\(\sum_iE_{i1}\otimes E_{1i}\) has product \(1\) and commutes with
the left and right action of every matrix unit. Its off-diagonal
bimodules are rectangular matrices times the respective arrows \(a,b\).
Their tensor products over the intermediate matrix block multiply
the rectangular matrices. This balanced multiplication is an isomorphism
on the corresponding matrix bimodules; an inverse inserts the selected
matrix unit through the first coordinate.

An \(H\)-balanced normalized cochain is therefore determined by the same
two coefficients as an \(E\)-balanced cochain of \(A_f\).
Explicitly, for alternating rectangular matrices \(M_1,\ldots,M_n\)
times their arrows, take the value of the corner cochain on those
arrows and multiply its scalar coefficient by \(M_1\cdots M_n\).
Use the resulting vertex or arrow block for the output.
In degree zero send the two scalars to the identities of the two
matrix blocks. Matrix units show that these formulas are inverse to
restriction to the selected corner.

This identification preserves the differential and every cup and
brace operation. An interior bar face has its product in \(H\), as
before. An endpoint multiplies the same matrices and arrow values.
For a brace, multiplying the matrices inside a slot and then outside
it gives the same matrix product by associativity. There is no trace
or dimension factor. The normalized \(H\)-bar is a projective diagonal
resolution: its terms split into finite sums of
\(\Lambda_de_i\otimes_R e_j\Lambda_d\), and the same extra-unit
contraction proves exactness. Thus the identification computes the
relative higher centre of \(\Lambda_d\), not only that of its corner.
It applies to the local parameter \(f=t\) and to the global parameter
\(f=s^2-\hbar^2\) in \eqref{eq:fc-global-order}.

For cochains of arities \(n,m\), define their Gerstenhaber bracket by
\([F,G]=F\{G\}-(-1)^{(n-1)(m-1)}G\{F\}\).
It has cochain degree \(-1\). The formulas just proved give
\begin{equation}\label{eq:fc-bracket}
 [\theta u^m,\theta u^n]=(m-n)\theta u^{m+n},\qquad
 [\theta u^m,u^n]=-n u^{m+n},\qquad [u^m,u^n]=0.
\end{equation}
Shifting the complex by \([1]\) makes this a degree-zero Lie bracket.
For the full deformation cochains the differential is
\([\mu,-]\), where \(\mu\) is multiplication; on arity \(n\) it is
\((-1)^{n-1}\) times the Hochschild differential above. On
\(\mathcal C_f\) the differentials agree, because the even cochains
are closed. The inclusion also compares the deformation Lie complexes:
use \(-h\) instead of \(h\) on odd arities of the full complex.

\section{A formal product with a nonformal bracket}

The complex \eqref{eq:fc-small} is a direct sum of \(R\) in degree
zero and the two-term free complexes
\(R\theta u^m\xrightarrow{f}Ru^{m+1}\). Hence
\begin{gather}\label{eq:fc-cohomology}
 H^0\mathcal C_f=R,\qquad
 H^{2m}\mathcal C_f=(R/fR)u^m\quad(m\geq1),\\
 H^{2m+1}\mathcal C_f=\operatorname{ann}_R(f)\theta u^m
 \quad(m\geq0).\notag
\end{gather}
The notation \(\operatorname{ann}_R(f)\) denotes the scalars killed
by multiplication by \(f\). The direct sum of these bounded free
complexes is \(K\)-flat: tensoring it with an acyclic complex remains
acyclic, term by term and then by direct sum. Its derived scalar change
therefore retains exactly the same generators, differential and braces.

\begin{theorem}\label{thm:fc-nonformal}
Suppose \(R\) is a commutative \(\mathbb Q\)-algebra and \(f\) is
a non-zero-divisor that is not a unit. The cup algebra
\(\mathcal C_f\) is formal over \(R\). Its deformation Lie algebra
\(\mathcal C_f[1]\) is not formal over \(R\), and cannot be joined
by Lie quasi-isomorphisms to any complex with zero bracket.
\end{theorem}
\begin{proof}
Regularity of \(f\) makes \(\theta\mapsto0\), \(u\mapsto u\)
a multiplicative quasi-isomorphism to \(H=R[u]/(fu)\).
This is the asserted cup formality. The cohomology is entirely even,
so its degree-minus-one bracket is zero.

Put \(B=R\oplus R\varepsilon\), with \(|\varepsilon|=-1\),
\(\varepsilon^2=0\), and \(d\varepsilon=f\). This commutative
differential graded algebra resolves \(R/fR\). Its underlying complex
is bounded and free, so tensoring any quasi-isomorphism with \(B\)
preserves it. The product on \(B\), with tensor signs, defines scalar
extension of a differential graded Lie algebra.

If a Lie quasi-isomorphism zigzag joined \(\mathcal C_f[1]\) to an
abelian differential graded Lie algebra, tensor that zigzag with
\(B\). The abelian end would still have zero cohomology bracket.
At the other end, \(K\)-flatness gives a Lie quasi-isomorphism
\(\mathcal C_f[1]\otimes B\to\mathcal C_0[1]\) over \(R/fR\).
The target differential is zero and
\([\theta,u]=-u\ne0\), since \(R/fR\ne0\).
A Lie quasi-isomorphism preserves the cohomology bracket, giving a
contradiction. This also rules out formality to \(H[1]\).
\end{proof}

Before specialization, the detecting cycles are
\(\Theta=\theta-\varepsilon u\) and \(U=u\) in
\(\mathcal C_f\otimes B\). They satisfy
\([\Theta,U]=-U\). Thus the lost operation has an explicit
representative. It is not inferred from a generic dimension count.
When \(f\) is invertible all positive degrees contract by
\(u^{m+1}\mapsto f^{-1}\theta u^m\); at \(f=0\), odd classes
survive. Ordinary specialization of the cohomology in
\eqref{eq:fc-cohomology} would miss those odd classes.

\section{A quadratic cover and two cochain representatives}

Assume now that \(2\) is invertible in \(R\). Let
\(S_f=R[z]/(z^2-f)\), and let \(g\) act by \(g(z)=-z\).
The skew group algebra \(S_f\rtimes\langle g\rangle\) consists
of \(r_0+r_1z+r_2g+r_3zg\), with \(g^2=1\) and
\(gz=-zg\). Multiplication uses these equations and \(z^2=f\).
There is an exact algebra isomorphism
\begin{equation}\label{eq:fc-root}
 A_f\simeq S_f\rtimes\langle g\rangle,\qquad
 e_0=\frac{1+g}{2},\quad e_1=\frac{1-g}{2},\quad
 a=\frac{z-zg}{2},\quad b=\frac{z+zg}{2}.
\end{equation}
Indeed \(z=a+b\), \(g=e_0-e_1\) are inverse formulas, and their
squares and anticommutation follow from \eqref{eq:fc-rankfour}.
This proves the isomorphism over zero divisors as well as over fields.

A module over this algebra is an \(S_f\)-module with a semilinear
involution, meaning an additive involution \(G\) with
\(G(sn)=g(s)G(n)\). Its eigenspaces split as \(N_+\oplus N_-\).
Multiplication by \(z\) gives maps \(a:N_-\to N_+\),
\(b:N_+\to N_-\), whose two composites are \(f\).
Conversely these two maps define the module. The same formulas apply
to morphisms and complexes. Thus the algebra remembers both characters
of the quadratic cover, even above its branch divisor.

The equivariant free line \(S_f\) has positive basis \(1\) and
negative basis \(z\); its \(z\)-matrix is
\(\left(\begin{smallmatrix}0&f\\1&0\end{smallmatrix}\right)\).
Changing the equivariance by the sign character exchanges these bases
and gives \(\left(\begin{smallmatrix}0&1\\f&0\end{smallmatrix}\right)\).
These are exactly the two modules \eqref{eq:fc-probes}. They are
geometric equivariant lines, not two choices of a faithful collision
matrix representation.

The cover also gives a symmetric choice of cochains. Put
\(U_m=(1,1)\) in even degree and \(V_m=(1/2,1/2)\) in odd degree
of the same two-coordinate complex. They span a second subcomplex
\(\mathcal D_f\), with
\begin{gather}\label{eq:fc-geometric-model}
 dU_m=0,\quad dV_m=fU_{m+1},\quad U_mU_n=U_{m+n},\\
 U_mV_n=V_nU_m=V_{m+n},\qquad
 V_mV_n=\frac f4U_{m+n+1}.\notag
\end{gather}
Its braces with only odd inner entries are
\[
 U_m\{V_{n_1},\ldots,V_{n_l}\}
 =2^{-l}\binom{2m}{l}U_{m+\sum n_i},\qquad
 V_m\{V_{n_1},\ldots,V_{n_l}\}
 =2^{-l}\binom{2m+1}{l}V_{m+\sum n_i}.
\]
An even inner entry again makes the brace zero. Every slot now
contributes, each with factor \(1/2\); this proves the formulas
directly. Projection onto the equal-coordinate subspace is a chain
map. Its complementary difference subspace is contracted by the
even-to-odd map with coefficient \(2\) in \eqref{eq:fc-periodic}.
Thus both inclusions in
\[
 \mathcal C_f\longrightarrow C^\bullet
 \longleftarrow\mathcal D_f
\]
are quasi-isomorphisms preserving the cup product and all braces.
Their endpoints need not have identical strict formulas:
\[
 V_0-\theta=d(e_1/2),\qquad V_0^2=fu/4,\qquad\theta^2=0.
\]
The averaging projection cannot be used as a multiplicative map:
it kills \(g=e_0-e_1\) but sends \(g^2=1\) to \(1\).
The inclusions provide the required comparison without making that
incorrect quotient.

The singular exchange has now led to an actual algebra, a resolution
of its relations, and its complete cochain insertions. The next
construction adds a cotangent variable to the quadratic cover. It
will preserve the two fusion modules while changing the geometric
category in a precisely specified way.
